% ERS_SRS_2B_v2.0.tex
% Especificacion de Requisitos de Software, AgroMoreira (SGA).
% Estandar base: ISO/IEC/IEEE 29148:2018 y ISO/IEC 25010:2023.
% Enfoque metodologico: legal-first (Amaral et al., 2021).
% Compilar: pdflatex -> bibtex -> pdflatex -> pdflatex

\documentclass[11pt,a4paper]{article}

\usepackage[utf8]{inputenc}
\usepackage[T1]{fontenc}
\usepackage[spanish,es-noshorthands]{babel}
\usepackage[a4paper,margin=2.4cm]{geometry}
\usepackage{longtable}
\usepackage{array}
\usepackage{booktabs}
\usepackage{enumitem}
\usepackage{fancyhdr}
\usepackage{titlesec}
\usepackage[hidelinks]{hyperref}
\usepackage{parskip}

\newcolumntype{L}[1]{>{\raggedright\arraybackslash}p{#1}}

\setlist{nosep,leftmargin=1.4em}
\titleformat{\section}{\large\bfseries}{\thesection}{0.6em}{}
\titleformat{\subsection}{\normalsize\bfseries}{\thesubsection}{0.5em}{}

\newcommand{\ruta}[1]{\texttt{#1}}
\newcommand{\rf}[2]{\subsection*{#1. #2}\addcontentsline{toc}{subsection}{#1. #2}}

\pagestyle{fancy}
\fancyhf{}
\fancyhead[L]{\small ERS/SRS AgroMoreira}
\fancyhead[R]{\small Version 2.0}
\fancyfoot[C]{\small Pagina \thepage}
\renewcommand{\headrulewidth}{0.4pt}

\begin{document}

\begin{titlepage}
\centering
\vspace*{2cm}
{\LARGE\bfseries Especificacion de Requisitos de Software\par}
\vspace{0.4cm}
{\large (ERS / SRS)\par}
\vspace{1.2cm}
{\Large\bfseries AgroMoreira\par}
{\large Sistema de Gestion Agricola con Inteligencia Artificial\par}
\vspace{2cm}
\begin{tabular}{rl}
\textbf{Version:} & 2.0 \\[2pt]
\textbf{Estandar base:} & ISO/IEC/IEEE 29148:2018 y ISO/IEC 25010:2023 \\[2pt]
\textbf{PFC:} & \#11, Categoria C, riesgo minimo operativo \\[2pt]
\textbf{Paralelo:} & 4to Software A \\[2pt]
\textbf{Fecha:} & 1 de septiembre de 2026 \\
\end{tabular}
\vspace{1.6cm}

\begin{tabular}{rl}
\textbf{Equipo:} & Jeanpierre Robinson Espinoza, lider \\
 & Danela Dayana Arteaga Alava, secretaria \\
 & Kamila Annabella Calle Delgado, documentadora \\
 & Maria del Rosario Escudero Plaza, modeladora \\
 & Roselyn Andreina Sanchez Centeno, verificadora \\[6pt]
\textbf{Docente supervisor:} & Ing. Gleiston Ciceron Guerrero Ulloa, PhD \\
\end{tabular}
\vfill
{\small Universidad Tecnica Estatal de Quevedo \textperiodcentered{} Facultad de Ciencias de la Computacion \textperiodcentered{} Ingenieria de Requerimientos (ISR-401)}
\end{titlepage}

\tableofcontents
\newpage

\section*{Historial de versiones}
\addcontentsline{toc}{section}{Historial de versiones}

\begin{longtable}{L{1.2cm} L{2.4cm} L{9.5cm} L{2.6cm}}
\toprule
\textbf{Version} & \textbf{Fecha} & \textbf{Cambios} & \textbf{Autor} \\
\midrule
\endfirsthead
\toprule
\textbf{Version} & \textbf{Fecha} & \textbf{Cambios} & \textbf{Autor} \\
\midrule
\endhead
1.0 & 30 de agosto de 2026 & Esqueleto inicial del documento, sin trabajo de campo todavia. & Equipo AgroMoreira \\
\addlinespace
1.1 & 31 de agosto de 2026 & Correccion del equipo y del enfoque metodologico a partir del expediente etico real. & Equipo AgroMoreira \\
\addlinespace
1.2 & 1 de septiembre de 2026 & Correccion del liderazgo del equipo. & Equipo AgroMoreira \\
\addlinespace
1.3 & 1 de septiembre de 2026 & Definicion del enfoque legal-first como enfoque oficial del proyecto. & Equipo AgroMoreira \\
\addlinespace
2.0 & 1 de septiembre de 2026 & Cierre de la Seccion 3 con los 39 requisitos funcionales y los 15 no funcionales obtenidos en las 16 entrevistas y las 9 sesiones de validacion, y de las historias de usuario con criterios de aceptacion para los 17 requisitos Must have. Cierre del modelado del sistema en la Seccion 4. Matriz de trazabilidad cerrada en 60 filas y Seccion 6 con el producto minimo viable funcional. & Equipo AgroMoreira \\
\bottomrule
\end{longtable}

\newpage

\section{Introduccion}

\subsection{Proposito}

Este documento especifica los requisitos funcionales y no funcionales del Sistema de Gestion Agricola con Inteligencia Artificial para Agricola Moreira, finca de policultivo dedicada al cacao y al platano verde. El componente empirico del proyecto sigue el enfoque legal-first de Amaral, Abualhaija, Sabetzadeh y Briand \cite{amaral2021}. A partir de un modelo conceptual de cumplimiento legal organizado en tres bloques normativos, la Ley Organica de Proteccion de Datos Personales del Ecuador, la Resolucion 183 de AGROCALIDAD sobre buenas practicas y trazabilidad del cacao, y la Resolucion 0072 sobre bioseguridad y manejo fitosanitario, se definieron 26 criterios de cumplimiento documentados en \ruta{Modelo\_Legal\_LOPDP.md}. El trabajo de campo permitio evaluar cuales de esos criterios quedaban cubiertos por los requisitos obtenidos mediante entrevistas convencionales y cuales solo se cubren cuando se derivan requisitos directamente del texto legal.

\subsection{Alcance}

El sistema cubre la centralizacion del registro de parcelas, cultivos y variedades, el registro de actividades agricolas y cosechas, la gestion del inventario de insumos, la asignacion y el seguimiento de tareas, la generacion de reportes de produccion y costos, y un modulo de alertas de plagas asistido por inteligencia artificial que siempre requiere confirmacion humana antes de aplicarse a un lote. Tambien cubre los requisitos de cumplimiento legal derivados del modelo de proteccion de datos y de la normativa agroindustrial ecuatoriana. El alcance definitivo se fijo tras las 16 entrevistas y las 9 sesiones de validacion descritas en la Seccion 3, y ya fue confirmado con el administrador Raul Moreira Alay.

\subsection{Glosario de terminos}

\begin{longtable}{L{3.2cm} L{12cm}}
\toprule
\textbf{Termino} & \textbf{Definicion} \\
\midrule
\endhead
RF & Requisito funcional. \\
\addlinespace
RNF & Requisito no funcional. \\
\addlinespace
RD & Requisito derivado directamente del texto legal, sin evidencia de entrevista. \\
\addlinespace
Sugerencia explicable & Recomendacion generada por el componente de inteligencia artificial, acompanada de una justificacion comprensible para el usuario y de los datos en los que se basa, sin que el sistema ejecute la accion de forma autonoma. \\
\addlinespace
AGROCALIDAD & Agencia de Regulacion y Control Fito y Zoosanitario del Ecuador. \\
\addlinespace
LOPDP & Ley Organica de Proteccion de Datos Personales del Ecuador. \\
\addlinespace
Perfil tecnico & Usuario con familiaridad previa en el manejo de sistemas informaticos, como el personal de administracion o los tecnicos agronomos. \\
\addlinespace
Perfil no tecnico & Usuario sin familiaridad previa en sistemas informaticos, como buena parte del personal de campo. \\
\bottomrule
\end{longtable}

\subsection{Referencias}

Las fuentes citadas en este documento estan en \ruta{referencias.bib}, con verificacion cruzada de autor, titulo, publicacion, ano y DOI. Esta bibliografia no sustituye a la del protocolo de investigacion.

\subsection{Vision general del documento}

La Seccion 2 describe el producto y sus interesados. La Seccion 3 presenta los 39 requisitos funcionales, los 15 requisitos no funcionales y las historias de usuario con sus criterios de aceptacion, todos obtenidos del trabajo de campo real. La Seccion 4 recoge el modelado del sistema, con los casos de uso, los diagramas de clases, de comportamiento y de componentes, y los mockups. Las Secciones 5 y 6 describen la priorizacion, la trazabilidad cerrada y el producto minimo viable.

\section{Descripcion general}

\subsection{Perspectiva del producto}

AgroMoreira es un sistema nuevo e independiente, desarrollado especificamente para Agricola Moreira. No reemplaza un sistema legado formal, ya que la finca lleva sus registros actuales en papel y en hojas de calculo, segun se confirmo durante las entrevistas.

\subsection{Diagrama de contexto}

\begin{verbatim}
                +---------------------------------------+
 Propietario /  |                                       |   Tecnicos de
 Administrador -+                                       +-- AGROCALIDAD
 de la finca    |                                       |   (externo)
                |                                       |
 Ingeniero      |         SISTEMA AGROMOREIRA           |   Centro de acopio /
 agronomo    ---+   (gestion agricola + sugerencias     +-- Asociacion de
 residente      |    de IA explicables)                 |   productores (externo)
                |                                       |
 Capataces   ---+                                       |
 Personal de    |                                       |
 campo       ---+                                       |
 Personal de    |                                       |
 acopio      ---+                                       |
                +---------------------------------------+
\end{verbatim}

La version trazada de este diagrama es \ruta{03\_Modelado/Diagrams/CTX01\_Context\_Diagram.drawio}, y la matriz de poder e interes es \ruta{CTX02\_Power\_Interest\_Matrix.drawio}.

\subsection{Mapa de interesados, matriz de poder e interes}

\begin{longtable}{L{3.4cm} L{1.4cm} L{1.4cm} L{2.6cm} L{5.4cm}}
\toprule
\textbf{Interesado} & \textbf{Poder} & \textbf{Interes} & \textbf{Cuadrante} & \textbf{Rol en el proyecto} \\
\midrule
\endhead
Propietario o administrador de la finca & Alto & Alto & Gestionar de cerca & Decide la adopcion del sistema y firmo el aval institucional. \\
\addlinespace
Ingeniero agronomo residente & Medio & Alto & Mantener informado e involucrar & Informante clave tecnico y usuario del modulo de IA. \\
\addlinespace
Capataces & Medio & Medio & Mantener satisfecho & Coordina al personal de campo y usa el sistema como usuario intermedio. \\
\addlinespace
Personal de campo & Bajo & Alto & Mantener informado & Usuario final principal de las sugerencias explicables. \\
\addlinespace
Personal de acopio & Bajo & Medio & Monitorear & Usuario del modulo de trazabilidad y cosecha. \\
\addlinespace
Tecnicos de AGROCALIDAD, externo & Medio & Bajo & Monitorear & Informante externo sobre requisitos normativos. \\
\addlinespace
Centro de acopio o asociacion, externo & Bajo & Bajo & Monitorear & Informante externo sobre trazabilidad posterior a la cosecha. \\
\bottomrule
\end{longtable}

\subsection{Modelado organizacional i estrella}

El propietario o administrador depende del personal de campo para la ejecucion correcta de las labores culturales. El personal de campo depende del sistema AgroMoreira para recibir sugerencias comprensibles y accionables sobre el manejo del cultivo. El propio sistema depende del ingeniero agronomo para validar la calidad de esas sugerencias de inteligencia artificial. El propietario o administrador depende tambien del sistema para la trazabilidad y los reportes que debe presentar ante AGROCALIDAD. El personal de acopio depende del sistema para registrar la cosecha y la calidad posterior a esta.

Los diagramas de dependencia estrategica y de razon estrategica estan en \ruta{03\_Modelado/Diagrams/ISTAR01\_Strategic\_Dependency\_SD.drawio} y \ruta{ISTAR02\_Strategic\_Rationale\_SR.drawio}.

\subsection{Entorno operativo}

El personal de campo trabaja principalmente con telefonos Android de gama baja y con poca experiencia previa en aplicaciones moviles, segun se confirmo en las entrevistas a los jornaleros. La conectividad dentro de las parcelas es intermitente, lo que llevo a definir el requisito no funcional de robustez de campo de la Seccion 3.3. El personal tecnico y administrativo cuenta con acceso a computadores para las tareas de configuracion, reporte y analisis.

\subsection{Restricciones de diseno}

El sistema debe operar con conectividad intermitente o nula en el campo, tal como confirmo el personal de campo entrevistado. Las sugerencias de inteligencia artificial deben ser explicables y en ningun caso pueden ejecutar acciones de forma autonoma sobre un lote sin confirmacion humana. El sistema debe cumplir con la Ley Organica de Proteccion de Datos Personales del Ecuador \cite{lopdp2021}, segun se detalla a continuacion.

\subsection{Mapeo a la Ley Organica de Proteccion de Datos Personales del Ecuador}

Fuente legal verificada, Registro Oficial, Quinto Suplemento N.\textsuperscript{o} 459, del 26 de mayo de 2021.

\begin{longtable}{L{2.4cm} L{4cm} L{9cm}}
\toprule
\textbf{Articulo} & \textbf{Tema} & \textbf{Relevancia para AgroMoreira} \\
\midrule
\endhead
Art. 7 y 8 & Tratamiento legitimo y consentimiento & El consentimiento debe ser libre, especifico, informado e inequivoco, y revocable en cualquier momento. Es la base del proceso de entrevistas y del formulario de consentimiento informado. \\
\addlinespace
Art. 10, literal e & Pertinencia y minimizacion de datos & El sistema y el estudio solo recolectan los datos estrictamente necesarios para la gestion agricola. \\
\addlinespace
Art. 10, literal g & Confidencialidad & El tratamiento no puede comunicarse para un fin distinto al recogido. Respalda el compromiso de confidencialidad firmado por el equipo. \\
\addlinespace
Art. 10, literal j & Seguridad de datos personales & Respalda el cifrado AES-256 de la zona restringida de evidencias. \\
\addlinespace
Art. 13 a 19 & Derechos de acceso, rectificacion, eliminacion, oposicion y portabilidad & El sistema en su version productiva debe soportar estos cinco derechos, recogidos en el RF-23. \\
\addlinespace
Art. 20 & Derecho a no ser objeto de una decision basada unica o parcialmente en valoraciones automatizadas & Es el articulo mas relevante para AgroMoreira. El titular puede exigir una explicacion motivada sobre cualquier sugerencia del sistema, conocer los datos usados e impugnarla. Da base legal directa a los requisitos de explicabilidad del RF-09 y de la Seccion 3.3. \\
\bottomrule
\end{longtable}

\subsection{Supuestos y dependencias}

Se confirmo que Agricola Moreira cuenta con dispositivos moviles y computadores disponibles para el personal tecnico y administrativo. La colaboracion del personal de campo, inicialmente asumida como voluntaria, quedo confirmada con la realizacion de las 16 entrevistas y las 9 sesiones de validacion descritas en la Seccion 3.

\section{Requisitos especificos}

Esta seccion reune los 39 requisitos funcionales y los 15 requisitos no funcionales del sistema, obtenidos de las 16 entrevistas semiestructuradas y de las 9 sesiones de validacion con walkthrough, 4 con perfiles tecnicos y 5 con perfiles no tecnicos. De los 39 requisitos funcionales, 31 cuentan con evidencia directa de entrevista y 8 se derivaron directamente de los 26 criterios de cumplimiento legal descritos en \ruta{Modelo\_Legal\_LOPDP.md}, en sus tres bloques normativos.

Cada entrevista tiene asignado un codigo de evidencia. EV-01 corresponde a ENTR-01, administrador. EV-02 a ENTR-02, jornalero. EV-03 a ENTR-03, jornalera. EV-04 a ENTR-04, trabajador. EV-05 y EV-06 a ENTR-05 y ENTR-06, tecnicos. EV-07 a ENTR-07, tecnico, en sesion de walkthrough. EV-08 a ENTR-08, jornalera, en sesion de walkthrough. EV-09 a ENTR-09, jornalero. EV-10 a ENTR-10, tecnico, en sesion de walkthrough. EV-11 y EV-12 a ENTR-11 y ENTR-12, jornaleros, en sesion de walkthrough. EV-13 y EV-14 a ENTR-13 y ENTR-14, jornaleros. EV-15 a ENTR-15, jornalero, en sesion de walkthrough. EV-16 a ENTR-16, tecnica, en sesion de walkthrough. EV-17 a ENTR-17, tecnico. Las 16 personas entrevistadas son distintas entre si, aunque algunas comparten nombre de pila.

\subsection{Requisitos funcionales con evidencia de entrevista}

\rf{RF-01}{Registro de parcelas y lotes}
\begin{itemize}
\item \textbf{Descripcion:} el sistema debe permitir registrar y editar una parcela o lote con nombre o numero, sector, ubicacion, area, tipo de cultivo, variedad, cantidad de plantas y estado, activo o inactivo.
\item \textbf{Actor y origen:} Administrador, Tecnico, Jornalero. EV-01, EV-07, EV-09, EV-10, EV-11, EV-14, EV-15, EV-16.
\item \textbf{Entradas:} nombre, sector, ubicacion, area, cultivo, variedad, numero de plantas.
\item \textbf{Salidas:} ficha de parcela registrada.
\item \textbf{Precondicion:} usuario autenticado con permiso de edicion.
\item \textbf{Postcondicion:} parcela visible en el listado general.
\item \textbf{Prioridad MoSCoW:} Must have.
\item \textbf{Criterio de verificacion:} al crear una parcela con los 7 campos obligatorios, el sistema la muestra en el listado en menos de 2 segundos, sin campos vacios.
\item \textbf{Evidencia textual:} ``Saber exactamente el lote por su nombre y numero, que cultivo tiene y cuantas plantas han sido o estan en tratamiento'' (EV-15). ``Los datos deberian ser el codigo, el tamano, el cultivo y la ubicacion'' (EV-16).
\end{itemize}

\rf{RF-02}{Catalogo cerrado de cultivos y variedades}
\begin{itemize}
\item \textbf{Descripcion:} el sistema debe ofrecer una lista predefinida y editable por el administrador de cultivos, cacao, platano, palma y ampliable, y sus variedades, sin permitir texto libre del tipo ``otros''.
\item \textbf{Actor y origen:} Jornalero, Tecnico. EV-10, EV-11, EV-12.
\item \textbf{Entradas:} seleccion de cultivo o variedad desde el catalogo.
\item \textbf{Salidas:} cultivo o variedad asociado a la parcela.
\item \textbf{Precondicion:} catalogo previamente configurado.
\item \textbf{Postcondicion:} ningun registro de cultivo queda con valor ``otros'' sin especificar.
\item \textbf{Prioridad MoSCoW:} Should have.
\item \textbf{Criterio de verificacion:} intentar guardar un cultivo fuera del catalogo es rechazado por el sistema.
\item \textbf{Evidencia textual:} ``no que este platano, cacao ahi... Simplemente que permite escribir... No deberia estar asi'' (EV-11). ``agrupadas por secciones, cacao, platano, etcetera'' (EV-12).
\item \textbf{Nota de relacion:} este RF se amplia en RF-34, que exige que el catalogo soporte ademas un conjunto de atributos distinto por variedad, no solo por cultivo.
\end{itemize}

\rf{RF-03}{Registro de actividades agricolas}
\begin{itemize}
\item \textbf{Descripcion:} el sistema debe permitir registrar una actividad, fumigacion, fertilizacion, poda, limpieza o riego, indicando fecha, lote, producto o insumo usado y trabajador responsable.
\item \textbf{Actor y origen:} Jornalero, Tecnico. EV-03, EV-07, EV-08, EV-10, EV-14, EV-15, EV-16.
\item \textbf{Entradas:} fecha, lote, tipo de actividad, producto, trabajador.
\item \textbf{Salidas:} registro de actividad.
\item \textbf{Precondicion:} lote existente.
\item \textbf{Postcondicion:} actividad queda asociada al historial del lote.
\item \textbf{Prioridad MoSCoW:} Must have.
\item \textbf{Criterio de verificacion:} el sistema no permite guardar una actividad sin fecha, lote y trabajador.
\item \textbf{Evidencia textual:} ``que se registre lo que es la fecha, la parcela, la actividad, producto y el trabajador'' (EV-10). ``La informacion dispensable seria la actividad, la parcela, la fecha y tambien el responsable'' (EV-16).
\end{itemize}

\rf{RF-04}{Registro de cosecha por unidad especifica del cultivo}
\begin{itemize}
\item \textbf{Descripcion:} el sistema debe registrar la cosecha usando la unidad propia de cada cultivo, racimo o caja para platano, tacho o quintal para cacao, con conversion configurable a kilogramos o libras.
\item \textbf{Actor y origen:} Jornalero. EV-02, EV-08, EV-11, EV-12, EV-13, EV-14, EV-15.
\item \textbf{Entradas:} lote, fecha, cultivo, cantidad, unidad.
\item \textbf{Salidas:} registro de cosecha.
\item \textbf{Precondicion:} lote y cultivo existentes.
\item \textbf{Postcondicion:} cosecha sumada al acumulado del lote o periodo.
\item \textbf{Prioridad MoSCoW:} Must have.
\item \textbf{Criterio de verificacion:} al registrar cosecha de cacao el sistema exige unidad tacho o quintal, y de platano racimo o caja.
\item \textbf{Evidencia textual:} ``cuando es la cosecha de platano, es por racimo. Y cuando es el cacao, por tacho'' (EV-02). ``Me gustaria el total de la produccion y tambien lo que es del peso'' (EV-08).
\end{itemize}

\rf{RF-05}{Calculo automatico de rendimiento por lote y periodo}
\begin{itemize}
\item \textbf{Descripcion:} el sistema debe calcular automaticamente el promedio de produccion por lote y periodo, semanal o mensual, y permitir comparar entre periodos.
\item \textbf{Actor y origen:} Administrador, Jornalero, Tecnica. EV-01, EV-07, EV-11, EV-15, EV-16.
\item \textbf{Entradas:} historico de cosechas del lote.
\item \textbf{Salidas:} promedio y comparacion entre periodos.
\item \textbf{Precondicion:} existen 2 o mas registros de cosecha del lote.
\item \textbf{Postcondicion:} valor mostrado en el reporte del lote.
\item \textbf{Prioridad MoSCoW:} Must have.
\item \textbf{Criterio de verificacion:} el promedio calculado coincide con el calculo manual, suma dividida entre N, sobre un conjunto de prueba.
\item \textbf{Evidencia textual:} ``sacar un promedio de cuanto rinde cada lote'' (EV-15). ``Me gustaria que generara automaticamente lo que es el rendimiento y tambien los totales'' (EV-16).
\end{itemize}

\rf{RF-06}{Calculo automatico de ganancia neta}
\begin{itemize}
\item \textbf{Descripcion:} el sistema debe calcular ingresos menos egresos, insumos y mano de obra, y mostrar la ganancia neta por periodo.
\item \textbf{Actor y origen:} Tecnico, Jornalero. EV-10, EV-11.
\item \textbf{Entradas:} registros de ingresos y egresos.
\item \textbf{Salidas:} ganancia neta del periodo.
\item \textbf{Precondicion:} existen registros de ingreso y egreso en el periodo.
\item \textbf{Postcondicion:} valor visible en el modulo financiero.
\item \textbf{Prioridad MoSCoW:} Must have.
\item \textbf{Criterio de verificacion:} la ganancia es igual a la suma de ingresos menos la suma de egresos, verificado contra un caso de prueba manual.
\item \textbf{Evidencia textual:} ``de todo lo que se gasto se ingresa tambien todo lo que se vendio y menorando todo eso, lo que quedo de ganancia'' (EV-10).
\item \textbf{Nota de relacion:} ver tambien RF-26, que pide un modulo de registro manual de ingresos y egresos independiente de este calculo automatico.
\end{itemize}

\rf{RF-07}{Gestion de inventario de insumos}
\begin{itemize}
\item \textbf{Descripcion:} el sistema debe permitir dar de alta insumos, registrar cantidad disponible o usada y visualizar el stock restante.
\item \textbf{Actor y origen:} Administrador, Tecnico, Jornalero. EV-01, EV-03, EV-05, EV-06, EV-09, EV-10, EV-11, EV-12, EV-16.
\item \textbf{Entradas:} nombre del insumo, cantidad, unidad.
\item \textbf{Salidas:} stock disponible actualizado.
\item \textbf{Precondicion:} ninguna.
\item \textbf{Postcondicion:} stock reducido tras cada uso registrado.
\item \textbf{Prioridad MoSCoW:} Must have.
\item \textbf{Criterio de verificacion:} tras registrar el uso de 26 unidades de un stock de 30, el sistema muestra 4 disponibles.
\item \textbf{Evidencia textual:} ``dice 4 de 30, es decir, lo que tenia era 30 y ahora le queda 4'' (EV-10). ``Sobre cada insumo, lo que es la cantidad, el nombre y ademas del stock limit'' (EV-16).
\end{itemize}

\rf{RF-08}{Alertas automaticas de bajo stock}
\begin{itemize}
\item \textbf{Descripcion:} el sistema debe enviar una alerta, notificacion, SMS o correo, cuando el stock de un insumo baje de un umbral configurable, anticipandose al calendario de fertilizacion o fumigacion.
\item \textbf{Actor y origen:} Administrador, Jornalero, Tecnica. EV-01, EV-03, EV-08, EV-09, EV-10, EV-11, EV-16.
\item \textbf{Entradas:} umbral configurado, stock actual.
\item \textbf{Salidas:} notificacion al responsable.
\item \textbf{Precondicion:} umbral definido para el insumo.
\item \textbf{Postcondicion:} notificacion enviada y registrada.
\item \textbf{Prioridad MoSCoW:} Must have.
\item \textbf{Criterio de verificacion:} al bajar del umbral, la notificacion se envia en menos de 5 minutos.
\item \textbf{Evidencia textual:} ``en siete dias esta estimado que debes abonar el cacao... Y se necesita comprar'' (EV-11). ``Me gustaria que me avisara por medio de una notificacion en el celular'' (EV-08).
\item \textbf{Nota de relacion:} ver RF-30 para el umbral configurable como campo propio del insumo.
\end{itemize}

\rf{RF-09}{Alerta de plagas o enfermedades asistida por IA con verificacion humana obligatoria}
\begin{itemize}
\item \textbf{Descripcion:} el sistema debe generar una alerta de posible plaga o enfermedad por lote, basada en los datos ingresados, marcada explicitamente como sugerencia a confirmar en campo, nunca como diagnostico definitivo automatico.
\item \textbf{Actor y origen:} Administrador, Tecnico, Jornalero. EV-01, EV-06, EV-07, EV-08, EV-09, EV-10, EV-15, EV-16, EV-17.
\item \textbf{Entradas:} datos del lote, historico de tratamientos.
\item \textbf{Salidas:} alerta con opcion de confirmar o descartar.
\item \textbf{Precondicion:} modelo de IA entrenado disponible.
\item \textbf{Postcondicion:} alerta queda registrada con el resultado de la verificacion humana.
\item \textbf{Prioridad MoSCoW:} Must have.
\item \textbf{Criterio de verificacion:} ninguna alerta se aplica automaticamente a un lote sin que el usuario responsable la confirme.
\item \textbf{Evidencia textual:} ``yo confirmaria yendo a visualizar el terreno'' (EV-10). ``Se deberia poder revisar antes de tomar una decision o aplicacion'' (EV-07, EV-08). ``Primeramente que eso venga de una fuente confiable o de ya experiencias realizadas'' (EV-17).
\end{itemize}

\rf{RF-10}{Diagnostico de plagas por imagen}
\begin{itemize}
\item \textbf{Descripcion:} el sistema debe permitir subir una foto de una hoja o fruto y recibir una sugerencia de plaga o enfermedad probable y tratamiento recomendado.
\item \textbf{Actor y origen:} Jornalero. EV-15.
\item \textbf{Entradas:} foto en JPG o PNG.
\item \textbf{Salidas:} plaga probable y recomendacion.
\item \textbf{Precondicion:} conexion a internet.
\item \textbf{Postcondicion:} resultado registrado en el historial del lote.
\item \textbf{Prioridad MoSCoW:} Could have.
\item \textbf{Criterio de verificacion:} el sistema responde con una sugerencia en menos de 10 segundos tras subir la foto.
\item \textbf{Evidencia textual:} ``una foto de una hoja o de un fruto enfermo que el sistema le atiene rapido y que plaga es, que remedio exacto'' (EV-15).
\end{itemize}

\rf{RF-11}{Asignacion y seguimiento de tareas}
\begin{itemize}
\item \textbf{Descripcion:} el sistema debe permitir asignar una tarea a un trabajador por lote y fecha, con estados pendiente, en proceso, bloqueado y completado, y filtro por estado.
\item \textbf{Actor y origen:} Administrador, Tecnico, Jornalero. EV-01, EV-09, EV-10, EV-12, EV-16.
\item \textbf{Entradas:} lote, trabajador, fecha, tipo de tarea.
\item \textbf{Salidas:} tarea visible con su estado.
\item \textbf{Precondicion:} trabajador registrado.
\item \textbf{Postcondicion:} tarea filtrable por los 4 estados.
\item \textbf{Prioridad MoSCoW:} Must have.
\item \textbf{Criterio de verificacion:} al filtrar por el estado pendiente, el sistema solo muestra tareas en ese estado.
\item \textbf{Evidencia textual:} ``que tengas esa opcion de buscar en esas cuatro secciones'' (EV-12).
\end{itemize}

\rf{RF-12}{Notificacion movil de tareas asignadas}
\begin{itemize}
\item \textbf{Descripcion:} el sistema debe notificar al trabajador, por push o SMS, el lote, la fecha y la tarea asignada.
\item \textbf{Actor y origen:} Tecnico, Jornalera. EV-07, EV-08, EV-10.
\item \textbf{Entradas:} asignacion de tarea.
\item \textbf{Salidas:} notificacion push o SMS.
\item \textbf{Precondicion:} trabajador con dispositivo registrado.
\item \textbf{Postcondicion:} notificacion entregada.
\item \textbf{Prioridad MoSCoW:} Should have.
\item \textbf{Criterio de verificacion:} la notificacion incluye lote, fecha y tipo de tarea en el mensaje.
\item \textbf{Evidencia textual:} ``que la notificacion... Le avise por lote, en tal lote, fecha tal, hay un control'' (EV-10).
\end{itemize}

\rf{RF-13}{Registro de motivo de rechazo o perdida de fruta}
\begin{itemize}
\item \textbf{Descripcion:} el sistema debe permitir registrar, por lote y fecha, la cantidad de fruta rechazada y su motivo, enfermedad, plaga o dano mecanico.
\item \textbf{Actor y origen:} Jornalero. EV-14.
\item \textbf{Entradas:} lote, fecha, cantidad rechazada, motivo.
\item \textbf{Salidas:} registro de rechazo.
\item \textbf{Precondicion:} cosecha registrada ese dia.
\item \textbf{Postcondicion:} rechazo restado del total aprovechable.
\item \textbf{Prioridad MoSCoW:} Should have.
\item \textbf{Criterio de verificacion:} el total de fruta aprovechable excluye automaticamente lo marcado como rechazado.
\item \textbf{Evidencia textual:} ``cuantos salieron rechazados, ya sea por una enfermedad, una plaga o los danos'' (EV-14).
\end{itemize}

\rf{RF-14}{Trazabilidad de lote de exportacion}
\begin{itemize}
\item \textbf{Descripcion:} el sistema debe permitir registrar, por lote de exportacion, el estado de calidad, primera o segunda, la variedad sembrada y el ciclo de enfunde, fecha y color de cinta, para trazar el producto desde el campo hasta el empaque.
\item \textbf{Actor y origen:} Jornalero. EV-13, EV-14.
\item \textbf{Entradas:} color de cinta, semana de enfunde, calidad de caja.
\item \textbf{Salidas:} ficha de trazabilidad del lote de exportacion.
\item \textbf{Precondicion:} lote marcado como exportacion.
\item \textbf{Postcondicion:} historial completo desde el enfunde hasta el empaque consultable.
\item \textbf{Prioridad MoSCoW:} Should have.
\item \textbf{Criterio de verificacion:} dado un numero de caja, el sistema recupera el lote, la fecha de enfunde y el color de cinta correspondiente.
\item \textbf{Evidencia textual:} ``con el dato que el enfundador da, con eso se lo maneja para el dia de cosecha'' (EV-13). ``que en las pantallas de la cosecha se debe ver clarito el color de la cinta'' (EV-14).
\end{itemize}

\rf{RF-15}{Consulta de precio de mercado y estimacion de ingreso}
\begin{itemize}
\item \textbf{Descripcion:} el sistema debe permitir consultar el precio de mercado vigente del cacao o del platano y estimar el ingreso proyectado de una cosecha.
\item \textbf{Actor y origen:} Jornalero. EV-11.
\item \textbf{Entradas:} cantidad cosechada, precio de mercado, manual o de fuente externa.
\item \textbf{Salidas:} ingreso estimado.
\item \textbf{Precondicion:} cantidad de cosecha registrada.
\item \textbf{Postcondicion:} estimacion mostrada al usuario.
\item \textbf{Prioridad MoSCoW:} Could have.
\item \textbf{Criterio de verificacion:} el ingreso estimado es igual a la cantidad multiplicada por el precio ingresado, validado con un caso de prueba.
\item \textbf{Evidencia textual:} ``que consultara el precio de la materia prima... Un calculo de cuanto mas o menos uno va a ganar'' (EV-11).
\end{itemize}

\rf{RF-16}{Registro de condicion climatica y estado del lote}
\begin{itemize}
\item \textbf{Descripcion:} el sistema debe permitir anotar, junto a cada visita o actividad del lote, el estado climatico, lluvia o sol, y del terreno, accesibilidad.
\item \textbf{Actor y origen:} Jornalero. EV-12, EV-15.
\item \textbf{Entradas:} condicion climatica, observacion de acceso.
\item \textbf{Salidas:} campo adicional en el registro de actividad o parcela.
\item \textbf{Precondicion:} ninguna.
\item \textbf{Postcondicion:} dato consultable en el historial del lote.
\item \textbf{Prioridad MoSCoW:} Could have.
\item \textbf{Criterio de verificacion:} el campo de clima aparece disponible en el formulario de registro de parcela o actividad.
\item \textbf{Evidencia textual:} ``Puede faltar un espacio para anotar lo que es el estado del clima'' (EV-15).
\end{itemize}

\rf{RF-17}{Canal de reporte rapido por chat o voz}
\begin{itemize}
\item \textbf{Descripcion:} el sistema debe ofrecer un boton de reporte rapido, texto corto o nota de voz, para que un jornalero reporte una novedad sin llenar formularios extensos.
\item \textbf{Actor y origen:} Jornalero. EV-15.
\item \textbf{Entradas:} nota de voz o texto corto.
\item \textbf{Salidas:} reporte enviado al responsable del lote.
\item \textbf{Precondicion:} usuario autenticado.
\item \textbf{Postcondicion:} reporte visible para el administrador o tecnico.
\item \textbf{Prioridad MoSCoW:} Could have.
\item \textbf{Criterio de verificacion:} el reporte se envia en un maximo de 2 toques desde la pantalla principal.
\item \textbf{Evidencia textual:} ``un chat o un boton rapido de voz para reportar novedades'' (EV-15).
\end{itemize}

\rf{RF-18}{Registro de riesgo y seguridad laboral por lote}
\begin{itemize}
\item \textbf{Descripcion:} el sistema debe permitir marcar un lote con una alerta de riesgo, fauna peligrosa o terreno inestable, visible antes de asignar tareas ahi, y sugerir el equipo de proteccion personal requerido.
\item \textbf{Actor y origen:} Jornalero, Tecnico. EV-11, EV-17.
\item \textbf{Entradas:} tipo de riesgo, lote, equipo de proteccion recomendado.
\item \textbf{Salidas:} alerta visible al asignar tarea en ese lote.
\item \textbf{Precondicion:} ninguna.
\item \textbf{Postcondicion:} advertencia mostrada al asignar personal.
\item \textbf{Prioridad MoSCoW:} Should have.
\item \textbf{Criterio de verificacion:} al asignar una tarea en un lote marcado con riesgo, el sistema muestra la advertencia antes de confirmar.
\item \textbf{Evidencia textual:} ``hay riesgo de que haya animales... Como serpientes'' (EV-11). ``se dice que usar, que trajes usar, incluso botas... Ya que al momento de presenciar estos peligros esten bien cuidados'' (EV-17).
\end{itemize}

\rf{RF-19}{Reportes de produccion, costos y perdidas con graficos}
\begin{itemize}
\item \textbf{Descripcion:} el sistema debe generar reportes de produccion, costos, perdidas por plaga y rendimiento por parcela o periodo, con graficos de barra y circulares codificados por color.
\item \textbf{Actor y origen:} Administrador, Tecnico. EV-01, EV-05, EV-06, EV-07, EV-10, EV-16.
\item \textbf{Entradas:} rango de fechas, lote o lotes.
\item \textbf{Salidas:} reporte visual exportable.
\item \textbf{Precondicion:} existen datos en el rango seleccionado.
\item \textbf{Postcondicion:} reporte generado y descargable.
\item \textbf{Prioridad MoSCoW:} Must have.
\item \textbf{Criterio de verificacion:} el reporte generado refleja exactamente la suma de los registros del rango de fechas seleccionado.
\item \textbf{Evidencia textual:} ``reporte por parcela, costos, perdidas por plagas o rendimiento por periodo'' (EV-01).
\end{itemize}

\rf{RF-20}{Gestion de usuarios y control de acceso}
\begin{itemize}
\item \textbf{Descripcion:} el sistema debe requerir usuario y contrasena para ingresar, con roles diferenciados, administrador, tecnico y jornalero.
\item \textbf{Actor y origen:} Tecnico. EV-07, EV-08, EV-10.
\item \textbf{Entradas:} usuario, contrasena.
\item \textbf{Salidas:} sesion autenticada con permisos segun rol.
\item \textbf{Precondicion:} usuario previamente registrado.
\item \textbf{Postcondicion:} acceso concedido solo a las funciones de su rol.
\item \textbf{Prioridad MoSCoW:} Must have.
\item \textbf{Criterio de verificacion:} un usuario con rol jornalero no puede acceder al modulo financiero.
\item \textbf{Evidencia textual:} ``han creado como un tipo formulario, donde le pide lo que es un usuario y una contrasena'' (EV-10). ``El administrador deberia tener el control total y los trabajadores solo lo necesario'' (EV-07, EV-08).
\end{itemize}

\rf{RF-21}{Registro de visitas tecnicas periodicas}
\begin{itemize}
\item \textbf{Descripcion:} el sistema debe permitir registrar cada visita tecnica, fecha, hallazgos e informe adjunto, con recordatorio de la siguiente visita en un ciclo de 15 dias.
\item \textbf{Actor y origen:} Tecnico. EV-10.
\item \textbf{Entradas:} fecha de visita, informe, proxima fecha estimada.
\item \textbf{Salidas:} historial de visitas tecnicas por lote o finca.
\item \textbf{Precondicion:} ninguna.
\item \textbf{Postcondicion:} recordatorio generado para la siguiente visita.
\item \textbf{Prioridad MoSCoW:} Should have.
\item \textbf{Criterio de verificacion:} 15 dias despues de una visita registrada, el sistema genera un recordatorio automatico.
\item \textbf{Evidencia textual:} ``un control de la visita tecnica de 15 dias y pone el otro control que fue a realizarla despues de 15 dias'' (EV-10).
\end{itemize}

\rf{RF-26}{Registro manual de ingresos y egresos}
\begin{itemize}
\item \textbf{Descripcion:} el sistema debe ofrecer un modulo dedicado para registrar manualmente cada ingreso y egreso, concepto, monto y fecha, independiente del calculo automatico de ganancia neta del RF-06.
\item \textbf{Actor y origen:} Tecnico, Jornalera. EV-07, EV-08.
\item \textbf{Entradas:} concepto, monto, fecha, tipo, ingreso o egreso.
\item \textbf{Salidas:} listado de movimientos.
\item \textbf{Precondicion:} usuario autenticado.
\item \textbf{Postcondicion:} movimiento reflejado en el calculo de ganancia neta.
\item \textbf{Prioridad MoSCoW:} Must have.
\item \textbf{Criterio de verificacion:} un ingreso o egreso registrado manualmente aparece reflejado en el reporte financiero del periodo correspondiente.
\item \textbf{Evidencia textual:} ``Que permitiera controlar los ingresos y egresos'' (EV-07). ``Me gustaria que tuviera un registro de gastos y de ganancias tambien'' (EV-08).
\end{itemize}

\rf{RF-27}{Vista de tareas pendientes por trabajador}
\begin{itemize}
\item \textbf{Descripcion:} el sistema debe mostrar a cada trabajador una vista dedicada de sus propias tareas pendientes.
\item \textbf{Actor y origen:} Jornalero. EV-09.
\item \textbf{Entradas:} usuario autenticado.
\item \textbf{Salidas:} listado de tareas pendientes propias.
\item \textbf{Precondicion:} existen tareas asignadas al usuario.
\item \textbf{Postcondicion:} ninguna.
\item \textbf{Prioridad MoSCoW:} Should have.
\item \textbf{Criterio de verificacion:} un trabajador con 3 tareas pendientes ve exactamente esas 3 al abrir la vista, sin tareas de otros lotes o trabajadores.
\item \textbf{Evidencia textual:} ``Existe alguna funcion que considere indispensable y que todavia no hayamos completado. Ver las tareas pendientes'' (EV-09).
\end{itemize}

\rf{RF-28}{Campo de observaciones en el registro de tareas}
\begin{itemize}
\item \textbf{Descripcion:} el sistema debe permitir adjuntar una nota u observacion libre a cada tarea registrada como completada.
\item \textbf{Actor y origen:} Jornalero. EV-09.
\item \textbf{Entradas:} texto libre.
\item \textbf{Salidas:} observacion asociada a la tarea.
\item \textbf{Precondicion:} tarea existente.
\item \textbf{Postcondicion:} observacion visible al consultar el historial de la tarea.
\item \textbf{Prioridad MoSCoW:} Could have.
\item \textbf{Criterio de verificacion:} la observacion ingresada aparece al consultar el detalle de la tarea.
\item \textbf{Evidencia textual:} ``hay alguna informacion que considere que falta. Si, las observaciones de las tareas'' (EV-09).
\end{itemize}

\rf{RF-29}{Historial completo por parcela}
\begin{itemize}
\item \textbf{Descripcion:} el sistema debe ofrecer una vista de historial por parcela que agregue actividades, tratamientos y cosechas a lo largo del tiempo, no solo el estado actual.
\item \textbf{Actor y origen:} Tecnica. EV-16.
\item \textbf{Entradas:} identificador de parcela.
\item \textbf{Salidas:} linea de tiempo de eventos de la parcela.
\item \textbf{Precondicion:} la parcela tiene al menos un evento registrado.
\item \textbf{Postcondicion:} ninguna.
\item \textbf{Prioridad MoSCoW:} Should have.
\item \textbf{Criterio de verificacion:} al consultar el historial de una parcela con 5 eventos registrados en distintas fechas, los 5 aparecen ordenados cronologicamente.
\item \textbf{Evidencia textual:} ``El historial de cada parcela'' (EV-16).
\end{itemize}

\rf{RF-30}{Umbral de stock minimo configurable por insumo}
\begin{itemize}
\item \textbf{Descripcion:} el sistema debe permitir definir, para cada insumo, un valor numerico de stock minimo que dispare la alerta de bajo stock del RF-08.
\item \textbf{Actor y origen:} Tecnica. EV-16.
\item \textbf{Entradas:} cantidad minima por insumo.
\item \textbf{Salidas:} umbral guardado en la ficha del insumo.
\item \textbf{Precondicion:} insumo ya registrado.
\item \textbf{Postcondicion:} el umbral queda disponible para el mecanismo de alertas.
\item \textbf{Prioridad MoSCoW:} Should have.
\item \textbf{Criterio de verificacion:} al bajar el stock del insumo por debajo del valor configurado, se dispara la alerta del RF-08.
\item \textbf{Evidencia textual:} ``Sobre cada insumo, lo que es la cantidad, el nombre y ademas del stock limit'' (EV-16).
\end{itemize}

\rf{RF-31}{Deteccion de valores atipicos al ingresar datos}
\begin{itemize}
\item \textbf{Descripcion:} el sistema debe comparar cada dato numerico recien ingresado, por ejemplo la cantidad cosechada, contra el promedio historico del mismo lote y avisar si el valor se aleja de forma significativa, antes de guardarlo.
\item \textbf{Actor y origen:} Tecnico. EV-17.
\item \textbf{Entradas:} valor ingresado, historico del lote.
\item \textbf{Salidas:} advertencia de valor atipico.
\item \textbf{Precondicion:} existe historico suficiente del lote, minimo 3 registros previos.
\item \textbf{Postcondicion:} el usuario confirma o corrige el dato antes de guardar.
\item \textbf{Prioridad MoSCoW:} Could have.
\item \textbf{Criterio de verificacion:} un valor que se desvia mas de 2 desviaciones estandar del promedio historico dispara la advertencia antes de guardar.
\item \textbf{Evidencia textual:} ``que al momento de meter un dato, una variable, me indique... Que este valor de aqui estuvo mal'' (EV-17).
\end{itemize}

\rf{RF-32}{Via de integracion futura con sensores de campo}
\begin{itemize}
\item \textbf{Descripcion:} el sistema debe exponer una interfaz capaz de recibir automaticamente lecturas de sensores de clima, suelo o nutrientes, como alternativa al ingreso manual, para una fase posterior del proyecto.
\item \textbf{Actor y origen:} Tecnico. EV-17.
\item \textbf{Entradas:} lectura de sensor, formato a definir en fase posterior.
\item \textbf{Salidas:} dato incorporado al historial del lote sin intervencion manual.
\item \textbf{Precondicion:} sensor compatible configurado.
\item \textbf{Postcondicion:} ninguna.
\item \textbf{Prioridad MoSCoW:} Won't have en esta version, queda documentado para trabajo futuro.
\item \textbf{Criterio de verificacion:} no aplica en esta version, queda como extension documentada.
\item \textbf{Evidencia textual:} ``Que ya la aplicacion que estan generando tenga sensores... Que tomen esos datos automaticamente'' (EV-17).
\end{itemize}

\rf{RF-33}{Bandeja de sugerencias para disputar una recomendacion de IA}
\begin{itemize}
\item \textbf{Descripcion:} el sistema debe ofrecer un mecanismo para que el usuario registre formalmente su desacuerdo con una recomendacion de IA, distinto del simple confirmar o descartar del RF-09, y que quede disponible como registro de retroalimentacion.
\item \textbf{Actor y origen:} Tecnico. EV-17.
\item \textbf{Entradas:} recomendacion en cuestion, comentario del usuario.
\item \textbf{Salidas:} registro de desacuerdo asociado a la recomendacion.
\item \textbf{Precondicion:} existe una recomendacion de IA previa.
\item \textbf{Postcondicion:} el desacuerdo queda visible para revision posterior del equipo tecnico.
\item \textbf{Prioridad MoSCoW:} Should have.
\item \textbf{Criterio de verificacion:} un desacuerdo registrado aparece en un listado consultable, con la fecha y el comentario del usuario.
\item \textbf{Evidencia textual:} ``yo obviamente que daria mi observacion en una bandejita de sugerencias'' (EV-17).
\end{itemize}

\rf{RF-34}{Esquema de atributos configurable por variedad de cultivo}
\begin{itemize}
\item \textbf{Descripcion:} el catalogo de cultivos del RF-02 debe permitir definir campos y variables distintas por variedad, no solo por tipo de cultivo, dado que el platano y las distintas variedades de cacao manejan datos de manejo agronomico diferentes.
\item \textbf{Actor y origen:} Tecnico. EV-17.
\item \textbf{Entradas:} variedad seleccionada.
\item \textbf{Salidas:} formulario con los campos correspondientes a esa variedad.
\item \textbf{Precondicion:} variedad dada de alta en el catalogo con su conjunto de atributos.
\item \textbf{Postcondicion:} ninguna.
\item \textbf{Prioridad MoSCoW:} Must have.
\item \textbf{Criterio de verificacion:} al seleccionar una variedad de cacao, el formulario muestra un conjunto de campos distinto al que se muestra al seleccionar platano.
\item \textbf{Evidencia textual:} ``Existen muchas diferencias, en el platano se registra lo que es... En el cacao se registra desde su mazorca... Hay muchos tipos de cacao, el nacional, el CN51'' (EV-17).
\end{itemize}

\rf{RF-35}{Registro de analisis de suelo previo a la siembra}
\begin{itemize}
\item \textbf{Descripcion:} el sistema debe permitir registrar, por parcela, el resultado de un analisis de suelo previo al establecimiento de un cultivo, indicando tipo de suelo y aptitud.
\item \textbf{Actor y origen:} Tecnico. EV-17.
\item \textbf{Entradas:} tipo de suelo, resultado de aptitud, fecha del analisis.
\item \textbf{Salidas:} ficha de analisis de suelo asociada a la parcela.
\item \textbf{Precondicion:} parcela registrada.
\item \textbf{Postcondicion:} dato consultable antes de aprobar la siembra.
\item \textbf{Prioridad MoSCoW:} Could have.
\item \textbf{Criterio de verificacion:} una parcela con analisis de suelo registrado muestra el resultado de aptitud al consultar su ficha.
\item \textbf{Evidencia textual:} ``una calicata, para ver o realizar un analisis de suelo... nos va a decir si es apto o no es apto sembrar en dicho terreno'' (EV-17).
\end{itemize}


\end{document}
